\documentclass[12pt,draftclsnofoot,onecolumn]{IEEEtran}

\usepackage{amsmath,amsfonts,amssymb,bm}
\usepackage{amsthm}
\usepackage{booktabs,array,multirow}
\usepackage{cite}
\usepackage{url}
\usepackage[hidelinks]{hyperref}
\usepackage{microtype}
\usepackage{tabularx}
\newtheorem{proposition}{Proposition}
\newtheorem{lemma}{Lemma}

\begin{document}
	
	\title{Supplementary Material for ''Auditable Complexity-Controlled Takagi--Sugeno Forecasting with Sequential PAC-Bayesian Risk Assessment for Engineering Time-Series Applications''}
	
	\author{Muneef Abdulkareem Farea Ahmed}
	
	\maketitle
	
	
	\section{Properness of the Augmented Hierarchical Prior}
	\label{sup:prior}
	
	Condition on $\mathcal G_0$, so all grids, prior-only scalers, antecedent maps, candidate-specific prior means, prior variances, and realized radius rule counts are fixed. The augmented hypothesis is $h=(M,v,s,a)$ with
	$M=(F,p,\gamma,r,K)$.
	
	For ridge, the family-conditional structural mass over $(p,\gamma)$ sums to one. For fixed-$K$ TSK, the family-conditional executable mass is
	
	\begin{equation}
		\sum_{p,\gamma}\pi_p(p)\pi_\gamma(\gamma)
		\sum_{K\in\mathcal K_{\mathrm{exec}}}\pi_K(K)\le 1 .
	\end{equation}
	
	For each radius-controlled family, the executable mass is
	
	\begin{align}
		&\sum_{p,\gamma}\pi_p(p)\pi_\gamma(\gamma)
		\sum_{r\in\mathcal R}\pi_r(r)\pi_K(K(r))
		\mathbf 1\{r\text{ eligible}\}
		\nonumber\\
		&\hspace{2cm}\le
		\sum_{p,\gamma}\pi_p(p)\pi_\gamma(\gamma)
		\sum_{r\in\mathcal R}\pi_r(r)
		=1 ,
	\end{align}
	
	because $0<\pi_K(K)\le1$.
	
	Any remaining probability mass in a family block is assigned to a dummy hypothesis used only to complete the probability measure. This dummy hypothesis is not considered a deployable forecasting candidate.
	
	Multiplication by the normalized family mass, localization-variant mass $\omega(v)$, and scale mass $\xi(s)$ therefore yields a proper probability measure.
	
	On the augmented space, the posterior
	
	\begin{equation}
		Q=\delta_{\widehat M}\delta_{\widehat v}\delta_{\widehat s}
		Q_{\widehat M,\widehat v,\widehat s,\widehat\rho}
	\end{equation}
	
	has exact divergence
	
	\begin{align}
		\mathrm{KL}(Q\|P)
		={}&-\log\pi(\widehat M)
		-\log\omega(\widehat v)
		-\log\xi(\widehat s)
		\nonumber\\
		&+
		\mathrm{KL}
		\left(
		Q_{\widehat M,\widehat v,\widehat s,\widehat\rho}
		\middle\|
		P_{\widehat M,\widehat v,\widehat s}
		\right).
	\end{align}
	
	The final term follows the diagonal-Gaussian expression given in the main paper. Posterior-ratio selection is absorbed by the simultaneous quantifier over $Q\ll P$; the temperature grid is charged separately through $\nu(\lambda)$.
	
	\begin{table}[htbp]
		\centering
		\caption{Predefined structural grids and prior components.}
		\label{tab:prior_components}
		\small
		\begin{tabular}{lll}
			\toprule
			Component & Candidate space & Prior assignment\\
			\midrule
			Lag order $p$ & predefined lag grid & normalized discrete prior\\
			Regularization $\gamma$ & predefined grid & normalized discrete prior\\
			Radius $r$ & predefined radius grid & normalized discrete prior\\
			Rule count $K$ & $\{2,3,4,6,8,12\}$ & normalized discrete prior\\
			Model family & Ridge/TSK families & family-level prior mass\\
			\bottomrule
		\end{tabular}
	\end{table}
	
	\section{Time-Uniform Bounded-Loss Proof}
	\label{sup:proof}
	
	For fixed $h$, conditional Hoeffding's lemma gives
	
	\begin{equation}
		\mathbb E
		\left[
		\exp\{\lambda(\mu_t(h)-\ell_t(h))-\lambda^2/8\}
		\mid\mathcal G_{t-1}
		\right]
		\le1 .
	\end{equation}
	
	Thus
	
	\begin{equation}
		L_t(h)=
		\exp
		\left\{
		\lambda
		\sum_{i=1}^{t}
		(\mu_i(h)-\ell_i(h))
		-t\lambda^2/8
		\right\}
	\end{equation}
	
	is a nonnegative supermartingale.
	
	Its $P$-mixture remains a nonnegative supermartingale because $P$ is $\mathcal G_0$-measurable. Ville's inequality, followed by the Donsker--Varadhan change of measure, yields simultaneously over $t$ and $Q_t\ll P$
	
	\begin{equation}
		\lambda t
		\left(
		R_t^{\mathrm{seq}}(Q_t)
		-\widehat R_t(Q_t)
		\right)
		\le
		\frac{t\lambda^2}{8}
		+
		\mathrm{KL}(Q_t\|P)
		+
		\log\frac1{\delta\nu(\lambda)} .
	\end{equation}
	
	Dividing by $\lambda t$ gives the proposition stated in the main paper. A union bound over the finite temperature grid produces the factor $\nu(\lambda)$.
	The result used in this work is therefore a specialization of an established sequential PAC-Bayesian framework rather than a new general concentration inequality.	
	
	\section{Complete Synthetic Data-Generating Processes}
	\label{sup:dgp}
	
	The implementation and accompanying configuration files specify the exact generators used in the experiments. The implemented processes are:
	
	\begin{align}
		\text{AR(2): }Y_t&=0.6Y_{t-1}-0.2Y_{t-2}+\varepsilon_t,
		&
		\varepsilon_t&\sim\mathcal N(0,0.2^2),
		\\
		\text{SETAR: }Y_t&=
		\begin{cases}
			0.65Y_{t-1}+\varepsilon_t,&Y_{t-1}\le0,\\
			-0.45Y_{t-1}+0.25Y_{t-2}+\varepsilon_t,&Y_{t-1}>0,
		\end{cases}
		\\
		\text{NARMA-10: }Y_{t+1}
		&=
		0.3Y_t+
		0.05Y_t\sum_{j=0}^{9}Y_{t-j}
		+1.5U_{t-9}U_t+0.1,
		\\
		\text{Mackey--Glass: }Y_{t+1}
		&=
		Y_t+
		\frac{0.2Y_{t-17}}
		{1+Y_{t-17}^{10}}
		-0.1Y_t+\varepsilon_t .
	\end{align}
	
	The GARCH process is zero-mean GARCH(1,1),
	
	\begin{align}
		h_t&=0.05+0.10Y_{t-1}^{2}+0.85h_{t-1},
		\\
		Y_t&=\sqrt{h_t}Z_t ,
	\end{align}
	
	where $Z_t$ is a standardized Student-$t_5$ innovation.
	
	The structural-break process changes after 60\% of the retained sample:
	
	\begin{equation}
		Y_t=
		\mu_j+\phi_j(Y_{t-1}-\mu_j)+\varepsilon_{t,j}.
	\end{equation}
	
	Each process retains 2,000 observations after a burn-in period of 500 samples. Exact seeds and configuration details are provided with the accompanying implementation artifacts.
	
	
	\section{Operational Gate and Recovery Audit}
	\label{sup:operation}
	
	The Tetouan development gate required:
	
	\begin{itemize}
		\item untruncated certificate $\le0.20$;
		\item certification target clipping $\le5\%$;
		\item validation RMSE no worse than the 12-step seasonal-naive fallback.
	\end{itemize}
	
	Candidate selection then minimized validation RMSE using the frozen tie-breaking order.
	
	The PJM gate was frozen after the Tetouan development result and before test opening. A candidate was eligible only if all predefined conditions were satisfied, including certificate threshold, clipping constraint, validation improvements, block-wise stability, and worst-block degradation limits.
	
	Among eligible candidates, minimum validation asymmetric cost was preferred, followed by RMSE, certificate, dimension, rule count, family, radius, ridge penalty, and lag. If no candidate qualified, the seasonal-naive fallback was deployed.
	
	The first execution stopped due to an infrastructure interruption after the AEP pretest decision had been written and before any test outcome had been opened. The restart used unchanged code, data, configuration, and recorded hashes. This event is reported as a computational recovery rather than as an uninterrupted execution.
	
	
	\section{Secondary Post-Hoc Statistical Benchmarks}
	\label{sup:posthoc}
	
	The SARIMA and ETS analysis was designed after the original PJM test outcomes were opened. It is not part of the frozen gate and does not modify any original certificate, decision, or holdout claim.
	
	Hyperparameters were selected using the original validation segment. Selected models were fitted using pretest observations and evaluated using rolling one-step forecasts.
	
	\input{posthoc_energy_benchmarks}
	
	
	\section{Actual Multi-Rule TSK Examples}
	\label{sup:rules}
	
	The rule analysis is reconstructed from the exact selected development candidates.
	
	The SETAR example is a five-rule radius TSK with $p=3$, $r=1.75$, and $\gamma=1$.
	
	The Tetouan example is the selected zone-1 four-rule radius TSK with $p=12$, $r=1.5$, and $\gamma=1$.
	
	Antecedents are fitted using prior data and posterior-center consequents using the bound set exactly as described in the main paper.
	
	\input{fuzzy_rule_examples}
	
	The complete machine-readable table contains one row per rule and lag, including antecedent centers, spreads, consequent coefficients, intercepts, prior support, and mean certification firing weights.
	
	The Tetouan example is retained despite its negative test outcome to prevent interpretability from being conflated with operational success.
	
	
	\section{Additional Statistical Details}
	
	The synthetic paired bootstrap resamples complete series rather than individual forecast origins and uses 20,000 draws.
	
	Wilcoxon signed-rank tests are two-sided with Holm correction.
	
	Moving-block bootstrap resamples forecast-error blocks within each energy trajectory and recomputes nonlinear RMSE and asymmetric-cost improvements.
	
	The PJM block length is 14 days with 7- and 28-day sensitivity analyses. The Tetouan block length is seven days at two-hour resolution.
	
	
	\section{Reproducibility Checklist}
	
	\begin{table}[htbp]
		\centering
		\begin{tabular}{ll}
			\toprule
			Item & Availability\\
			\midrule
			Random seeds & Provided\\
			Chronological data splits & Defined\\
			Configuration files & Provided\\
			Model-selection rules & Defined\\
			Hyperparameter grids & Reported\\
			Statistical procedures & Reported\\
			Artifact records & Provided\\
			\bottomrule
		\end{tabular}
	\end{table}
	
	
\end{document}